% ===== PRÉAMBULE CANONIQUE — à coller VERBATIM en tête de CHAQUE .tex =====
\documentclass[11pt,a4paper]{article}
\usepackage[utf8]{inputenc}\usepackage[T1]{fontenc}\usepackage{lmodern}
\usepackage[french]{babel}
\usepackage{amsmath,amssymb,mathtools}\usepackage{icomma}
\usepackage[version=4]{mhchem}          % \ce{...} : équations & formules
\usepackage{siunitx}\sisetup{locale=FR} % \qty{}{}, \si{}, \ang{}
\usepackage{chemfig}                    % \chemfig{...} : structures développées
\usepackage{xcolor}\usepackage{colortbl}
\usepackage{tikz}\usepackage{pgfplots}\pgfplotsset{compat=1.18}
\usepackage[most]{tcolorbox}\usepackage{enumitem}\usepackage{array,booktabs}
\usepackage[a4paper,margin=2.2cm]{geometry}
\usetikzlibrary{arrows.meta,calc,angles,quotes,positioning,patterns,backgrounds,shapes.geometric}
\definecolor{primary}{RGB}{222,49,99}\definecolor{primarydark}{RGB}{150,22,55}
\definecolor{defcol}{RGB}{0,114,178}\definecolor{methcol}{RGB}{52,92,148}
\definecolor{excol}{RGB}{0,135,90}\definecolor{attncol}{RGB}{200,40,55}
\tcbset{boxbase/.style={enhanced,breakable,boxrule=0.9pt,arc=2.5pt,
  left=3mm,right=3mm,top=2.3mm,bottom=2.3mm,fonttitle=\bfseries,coltitle=white}}
% --- boîtes TUTORIEL (noms EXACTS, ne pas renommer) ---
\newtcolorbox{cle}[1][]{boxbase,colback=defcol!6,colframe=defcol,
  title={\textsc{À retenir}\ifx\relax#1\relax\relax\else\textmd{ : \itshape #1}\fi}}
\newtcolorbox{methode}[1][]{boxbase,colback=methcol!7,colframe=methcol,sharp corners,
  title={\textsc{Méthode}\ifx\relax#1\relax\relax\else\textmd{ --- \itshape #1}\fi}}
\newtcolorbox{exemplebox}[1][]{boxbase,colback=excol!5,colframe=excol,
  title={\textsc{Exemple}\ifx\relax#1\relax\relax\else\textmd{ : \itshape #1}\fi}}
\newtcolorbox{piege}[1][]{boxbase,colback=attncol!6,colframe=attncol,
  title={\textsc{Piège}\ifx\relax#1\relax\relax\else\textmd{ : \itshape #1}\fi}}
\newtcolorbox{remarque}[1][]{boxbase,colback=black!4,colframe=black!45,
  title={\textsc{Remarque}\ifx\relax#1\relax\relax\else\textmd{ : \itshape #1}\fi}}
% --- boîtes EXERCICE / CORRIGÉ (noms EXACTS, auto-comptées) ---
\newtcolorbox[auto counter]{exo}[1][]{boxbase,colback=primary!4,colframe=primary,
  title={\textsc{Exercice}~\thetcbcounter\ifx\relax#1\relax\relax\else\textmd{ --- \itshape #1}\fi}}
\newtcolorbox[auto counter]{corrige}[1][]{boxbase,colback=excol!4,colframe=excol,
  title={\textsc{Corrigé}~\thetcbcounter\ifx\relax#1\relax\relax\else\textmd{ --- \itshape #1}\fi}}
\newcommand{\R}{\mathbb{R}}\newcommand{\N}{\mathbb{N}}\newcommand{\Z}{\mathbb{Z}}\newcommand{\ds}{\displaystyle}
% (un \usepackage{fancyhdr}+en-tête est possible ; il est ignoré côté web)
% =========================================================================
\begin{document}
\sisetup{per-mode=symbol}
\begin{corrige}[du fer à l'oxyde : masse $\to$ masse]
\begin{enumerate}[label=\alph*)]
  \item $n(\ce{Fe})=\dfrac{11,2}{56}=\qty{0.2}{\mole}.$
  \item $n(\ce{Fe2O3})=n(\ce{Fe})\times\dfrac{2}{4}=0,2\times0,5=\qty{0.1}{\mole}$, d'où
        $m=0,1\times160=\qty{16}{\gram}.$
  \item $n(\ce{O2})=n(\ce{Fe})\times\dfrac{3}{4}=0,2\times0,75=\qty{0.15}{\mole}$, d'où
        $V=0,15\times22,4=\qty{3.36}{\litre}.$
\end{enumerate}
\end{corrige}

\begin{corrige}[volumes de gaz : le raccourci]
Les volumes se traitent comme des moles (mêmes conditions).
\begin{enumerate}[label=\alph*)]
  \item $V(\ce{O2})=2\times\dfrac{5}{1}=\qty{10}{\litre}.$
  \item $V(\ce{CO2})=2\times\dfrac{3}{1}=\qty{6}{\litre}.$
\end{enumerate}
\end{corrige}

\begin{corrige}[précipitation en solution]
\begin{enumerate}[label=\alph*)]
  \item $n(\ce{AgNO3})=C\,V=0,5\times0,100=\qty{0.05}{\mole}.$
  \item Rapport $1:1$ : $n(\ce{AgCl})=\qty{0.05}{\mole}$, d'où
        $m=0,05\times143,5\approx\qty{7.18}{\gram}.$
\end{enumerate}
\end{corrige}

\begin{corrige}[construire un tableau d'avancement]
\begin{enumerate}[label=\alph*)]
  \item $n(\ce{SO2})=0,4-2\xi$ ; \; $n(\ce{O2})=0,3-\xi$ ; \; $n(\ce{SO3})=2\xi$.
  \item \ce{SO2} épuisé : $0,4-2\xi=0\Rightarrow \xi=\qty{0.2}{\mole}.$
  \item $n(\ce{O2})=0,3-0,2=\qty{0.1}{\mole}$ ; \; $n(\ce{SO3})=2\times0,2=\qty{0.4}{\mole}.$
        (Il reste \qty{0.1}{\mole} de \ce{O2} : ce réactif était en excès.)
\end{enumerate}
\end{corrige}

\begin{corrige}[acétylène du chalumeau]
\begin{enumerate}[label=\alph*)]
  \item $n(\ce{CaC2})=\dfrac{320}{64}=\qty{5}{\mole}.$
  \item Rapport $1:1$ : $n(\ce{C2H2})=\qty{5}{\mole}.$
  \item $V=n\,V_{\mathrm{m}}=5\times22,4=\qty{112}{\litre}.$
\end{enumerate}
\end{corrige}

\end{document}
